\section{Related work}

Counterfactual inference for determining causal effects  in observational studies has been studied extensively in statistics, economics, epidemiology and sociology \citep{morgan2014counterfactuals,robins2000marginal,rubin2011causal,chernozhukov2013inference} as well as in machine learning \citep{langford2011doubly,bottou2013counterfactual,swaminathan2015batch}.


Non-parametric methods do not attempt to model the relation between the context, intervention, and outcome. The methods include nearest-neighbor matching, propensity score matching, and propensity score re-weighting \citep{rosenbaum1983central, rosenbaum2002observational,austin2011introduction}.

Parametric methods, on the other hand, attempt to concretely model the relation between the context, intervention, and outcome. These methods include any type of regression including linear and logistic regression \citep{prentice1976use,gelman2006data}, random forests \citep{wager2015estimation}  and regression trees \citep{chipman2010bart}.

Doubly robust methods combine aspects of parametric and non-parametric methods, typically by using a propensity score weighted regression \citep{bang2005doubly,dudik2011doubly}. They are especially of use when the treatment assignment probability is known, as is the case for off-policy evaluation or learning from logged bandit data. Once the treatment assignment probability has to be estimated, as is the case in most observational studies, their efficacy might wane considerably \citep{kang2007demystifying}.

\citet{tian2014simple} presented one of the few methods that achieve balance by transforming or selecting covariates, modeling  interactions between treatment and covariates.
